% Prospective 2016 wording repairs, re-targeted to main.tex @ 55c0c5a358eea193e6d2293a2c814a5435487d14
% (first written against 8df7527c).
% Numbers are unchanged; every figure below is recomputed in PROSPECTIVE_RECOMPUTATION.json.

% --- P-1 (still present at 55c0c5a3 line 365): replace lines 364-365 ("...eight sensitive-recovery clauses capping added recovery at $+0.001$ nats.")
eight sensitive-recovery clauses capping each release's measured recovery at no more than $0.001$ nats
above that of $J$.

% --- P-1 (continued): append after line 392 ("...not restate them as a jointly certified family.")
Measured recovery beyond the published service nevertheless remains positive for both releases
($0.0034$--$0.0060$ nats for $Q$ on six of eight endpoints, unadjusted), so passing these clauses means
leaking measurably less than $J$, not adding nothing.

% --- P-3: append after line 372 ("...neither establishes the complete claim.")
$Q$'s task point estimates ($-0.00327$, $-0.00400$) are themselves beyond the margin; the clauses fail on
precision, and $D_{17}$'s by less than $10^{-4}$. This is an effect of about the margin's size that the
registered test did not resolve, not an absent effect.

% --- P-2 (OPTIONAL at 55c0c5a3: lines 386-392 already state that the bounds are per-clause, not a jointly
% certified family, which is correct). If a positive statement is wanted, table caption line 395, replace "Upper bounds are one-sided $97.5\%$," by:
Upper bounds are pointwise one-sided $97.5\%$ (not simultaneous; all eight sensitive bounds and both task
signs also hold under a Bonferroni correction over the $20$ primary endpoints, post hoc and descriptive),

% --- optional, "Q is worse on unweighted task loss" sentence (formerly lines 401-402): after "$Q$ is worse on unweighted task loss" add
(the person-weighted difference is unresolved)
